% Prova sintética no formato do ENEM (prompt 03, §55) — nada aqui é conteúdo do INEP.
% Página 1: capa, uma coluna. Páginas 2 e 3: duas colunas com calha estreita, como os cadernos.
%   p2 c1: LINGUAGENS · inglês Q1 Q2 · espanhol Q1 Q2 · Q06 (a última da área, fora da faixa de língua)
%   p2 c2: MATEMÁTICA · Q136 normal · Q137 com gráfico · Q138 começa (atravessa a PÁGINA)
%   p3 c1: Q138 termina · Q139 começa (atravessa a COLUNA)
%   p3 c2: Q139 termina · Q140 com alternativas em fórmula · Q141, a última
\documentclass[10pt]{article}
% Preâmbulo comum das fixtures do scan (D48). Fontes Latin Modern com T1, para que a camada de
% texto traga os acentos como caracteres, e babel em português, para que \chapter imprima
% "Capítulo" como um livro brasileiro imprime.
\usepackage[T1]{fontenc}
\usepackage[utf8]{inputenc}
\usepackage[brazil]{babel}
\usepackage{lmodern}
\usepackage{amsmath,amssymb}
\usepackage{tikz}
\usepackage[a4paper,margin=2.5cm]{geometry}
\usepackage{enumitem}
\setlist[enumerate,2]{label=\alph*)}
\makeatletter\@ifundefined{c@chapter}{\newcounter{exemplo}}{\newcounter{exemplo}[chapter]}\makeatother
\newenvironment{exemplo}{\par\medskip\refstepcounter{exemplo}\noindent\textbf{Exemplo \theexemplo.}\ }{\par\medskip}
\newcommand{\blocoexercicios}[1][Exercícios]{\section*{#1}}

\usepackage{multicol,fancyhdr,tikz}
\raggedcolumns
\geometry{margin=1.6cm,top=2.2cm,bottom=2.2cm}
\setlength{\columnsep}{0.35cm}
\setlength{\columnseprule}{0.4pt}
\pagestyle{fancy}\fancyhf{}
\fancyhead[L]{\small ENEM 2099 \textbullet{} 2º DIA \textbullet{} CADERNO 9 \textbullet{} AZUL}
\fancyfoot[L]{\small *AZUL9 7 4AB1*}
\fancyfoot[R]{\small Página \thepage}
\setlength{\parindent}{0pt}
\newcommand{\questao}[1]{\par\medskip{\bfseries QUESTÃO #1}\par\smallskip}
\newcommand{\alt}[2]{\par\makebox[1.3em][l]{\bfseries #1}#2}
\newcommand{\area}[1]{\par\medskip{\centering\bfseries #1\par}\smallskip}
\newcommand{\enunciado}{Leia o texto a seguir com atenção. Este enunciado sintético existe
para ocupar algumas linhas da coluna, como um texto-base de verdade ocuparia numa prova.}
\newcommand{\alternativas}{\alt{A}{primeira alternativa.}\alt{B}{segunda alternativa.}%
\alt{C}{terceira alternativa.}\alt{D}{quarta alternativa.}\alt{E}{quinta alternativa.}\par}
\begin{document}
\thispagestyle{empty}
\begin{center}
{\LARGE\bfseries EXAME NACIONAL DO ENSINO MÉDIO}\\[1cm]
{\Large PROVA SINTÉTICA — 2º DIA}\\[2cm]
\end{center}
Instruções: esta capa não contém questões e não pode entrar em recorte nenhum.
\newpage
\begin{multicols*}{2}
\area{LINGUAGENS, CÓDIGOS E SUAS TECNOLOGIAS}
\area{Questões de 01 a 05 (opção inglês)}
\questao{01}\enunciado\alternativas
\questao{02}\enunciado\alternativas
\area{Questões de 01 a 05 (opção espanhol)}
\questao{01}\enunciado\alternativas
\questao{02}\enunciado\alternativas
\questao{06}\enunciado\alternativas
\columnbreak
\area{MATEMÁTICA E SUAS TECNOLOGIAS}
\area{Questões de 136 a 180}
\questao{136}\enunciado\alternativas
\questao{137}Observe o gráfico.\par
\begin{center}\begin{tikzpicture}[scale=0.8]
\draw[->] (0,0) -- (4,0); \draw[->] (0,0) -- (0,3);
\draw[thick] (0,0.5) -- (1,2) -- (2,1.2) -- (3,2.6);
\end{tikzpicture}\end{center}
Qual é o maior valor mostrado?\alternativas
\questao{138}\enunciado\par \enunciado\par
Esta questão continua na próxima página.\par
\columnbreak
Continuação da questão 138, já na página seguinte.\alternativas
\questao{139}\enunciado\par \enunciado\par
Esta questão continua na próxima coluna.\par
\columnbreak
Continuação da questão 139, na coluna da direita.\alternativas
\questao{140}Qual expressão é igual a $\frac{1}{2}+\frac{1}{3}$?\par
\alt{A}{$\frac{5}{6}$}\alt{B}{$\frac{2}{5}$}\alt{C}{$\frac{1}{6}$}\alt{D}{$\frac{1}{5}$}\alt{E}{$\sqrt{2}$}\par
\questao{141}\enunciado\alternativas
\end{multicols*}
\end{document}
